\section{Towards Weakly Supervised Temporal Uncertainty Localization}
\label{sec:non_static}

Mechanistic detectors face a tension among coverage, localization,
and scalability. Fixed-position probes \citep{azaria2023,orgad2024}
are efficient but read only one layer-token slice, so they miss
signals that appear elsewhere in the activation tensor. Full-response
activation maps cover broader structure, but pooling the whole
response produces a sequence-level score and scales poorly for
localization \citep{barshalom2025}. Claim-level detectors provide
coarser localization, but scoring extracted claims in isolation
ignores boundary evidence from neighboring supported and hallucinated
spans \citep{snyder2024,he2024}. TPE therefore uses local activation
windows as the inference unit: each score is tied to a bounded token
neighborhood, while still covering evidence before, during, and after
the labeled span. Appendix Figure~\ref{fig:signal_propagation} shows
that hallucination signal appears before, within, and after
hallucinated spans.

\paragraph{Why encode activations as 2D images?}
Per-layer and per-token probes show that useful hidden-state signal
is not concentrated in one fixed position. Informative regions spread
across both the layer and token axes, and the most predictive
layer-token cells shift across LLMs, as shown by the heatmaps in
Appendix~\ref{app:heatmaps}. TPE preserves this local layer-token
plane instead of using one fixed slice or one-axis pooling. Following
the activation-image view of ACT-ViT \citep{barshalom2025actvit}, it
encodes each local activation window with Perception Encoder features
\citep{bolya2025perceptionencoder}.

\paragraph{Weakly supervised temporal localization.}
This representation turns hallucination detection into weakly
supervised temporal localization. As in video localization, where
clip-level labels supervise event discovery over time
\citep{pilhyeonlee2021}, a response label or claim label supervises
a sequence of candidate activation windows without window-level
annotations. Section~\ref{sec:model} formalizes this view: TPE
encodes local layer-token windows, propagates evidence with a
temporal head, and uses top-$k$ MIL so sparse hallucination evidence
can be assigned to a small subset of windows.
